\documentclass[11pt,a4paper]{article}

\usepackage[margin=1in]{geometry}
\usepackage{booktabs}
\usepackage{longtable}
\usepackage{array}
\usepackage{graphicx}
\usepackage{caption}
\usepackage{subcaption}
\usepackage{xurl}
\usepackage{hyperref}
\usepackage{amsmath}
\usepackage{amssymb}
\usepackage{setspace}
\usepackage{float}
\usepackage{pdflscape}
\usepackage{adjustbox}

\onehalfspacing

\title{Supplementary Material for GWASRanker}
\author{Muhammad Muneeb and David B. Ascher}
\date{}

\begin{document}

\maketitle

\section*{Supplementary Section 1: Overview of the GWASRanker feature dataset}

The GWASRanker supplementary analyses were based on a merged feature table in which each row represented one candidate GWAS summary-statistic file for one phenotype. The table combined file-level metadata, schema-mapping features, harmonisation outputs, sample-size summaries, allele and variant-quality features, reference-panel overlap metrics, target-genotype compatibility metrics, clumping and pruning features, target Fisher GWAS concordance features, and downstream PRS performance summaries.

The merged feature table contained 284 candidate GWAS files and 1,015 columns. These columns were organised into base metadata, processing-status indicators, Feature1--Feature9 outputs, PLINK PRS summaries, PRSice-2 summaries, and additional PRS performance fields. This supplementary material documents the major sources of variation captured across these feature groups and describes how they were used to support GWASRanker model development.

\subsection*{Supplementary Section 1.1: Feature-family composition}

Feature columns were grouped according to the module that generated them. Feature1 and Feature2 described file-level structure and raw column mapping. Feature3 and Feature5 described harmonisation, genome-build detection, standard-column availability, and GWASLab-derived summaries. Feature4 described sample-size, allele, variant identifier, P-value, INFO, and allele-frequency summaries. Feature6 described heritability-related, statistical-inflation, reference allele-frequency, HapMap3, and LDSC-related features. Feature7 described reference-panel overlap and clumping-related summaries. Feature8 described overlap between each GWAS and the target genotype data. Feature9 described effect-size distributions, PRS-readiness, target Fisher GWAS concordance, clumping, and LD-pruning outputs. PLINK and PRSice-2 columns described downstream PRS performance.

For each candidate GWAS file, processing-status columns were used to determine whether each feature module completed successfully. These status columns were retained in the raw master table to support auditing and troubleshooting, but were not used as predictive features in GWASRanker models unless explicitly converted into binary processing-completion indicators.

The raw merged table was filtered before downstream modelling. Path columns, file-location columns, log-file columns, redundant metadata columns, constant columns, all-missing columns, duplicate-content columns, and outcome-leakage columns were removed. PRS performance columns were retained only as outcome variables. The remaining non-constant numeric and interpretable categorical features were used to build GWASRanker Model 1 and Model 2.

\begin{table}[H]
\centering
\caption{\textbf{Summary of major GWASRanker feature groups.}}
\label{tab}
\begin{tabular}{p{0.18\textwidth}p{0.72\textwidth}}
\toprule
\textbf{Feature group} & \textbf{Main information captured} \\
\midrule
Feature1 & File extension, compression, reader, file size, delimiter, header detection, and physical line counts. \\
Feature2 & Raw headers, normalised headers, mapped standard columns, missingness, and PRS-minimum-column compatibility. \\
Feature3 & Genome-build detection, GWASLab loading, harmonised variant count, standard-column loading, and QC removal. \\
Feature4 & Sample size, allele representation, variant identifiers, P-value summaries, INFO summaries, and allele-frequency summaries. \\
Feature5 & FinalGWAS-level QC, chromosome representation, duplicate variants, allele quality, P-value validity, and PRS/LDSC/PLINK compatibility. \\
Feature6 & Heritability-related features, genomic inflation, beta and SE distributions, reference allele-frequency overlap, HapMap3 overlap, and LDSC outputs. \\
Feature7 & Reference overlap with HapMap3, dbSNP151, 1000 Genomes EUR, PLINK BIM references, and clumping-related outputs. \\
Feature8 & GWAS--target genotype overlap, allele compatibility, MAF concordance, target phenotype sample characteristics, and final usable variants. \\
Feature9 & Effect-size distributions, PRS-readiness, target Fisher GWAS concordance, clumping outputs, and LD-pruning outputs. \\
PRS outputs & PLINK and PRSice-2 train, test, null, pure PRS, best-model, and generalisation-gap performance summaries. \\
\bottomrule
\end{tabular}
\end{table}
\subsection*{Supplementary Section 1.2: File- and schema-level variation}

\subsection{File and schema variation}

File and schema variation was assessed using Feature1 and Feature2, which captured the technical structure of the raw GWAS summary-statistic files and their compatibility with the standard GWASRanker schema. Across the 284 candidate GWAS files, variation was observed at multiple levels: file format and compression,delimiter structure, file size and scale, column count, header schema diversity,standard-column mapping completeness, effect-size and P-value representation,allele encoding, per-file missingness and duplication rates, and PRS preprocessing requirements.

\paragraph{File-level structure.}The large majority of files used the \texttt{.tsv.gz} format (264 of 283 fileswith available extension information; 93.0\%), with smaller numbers using\texttt{.txt.gz} ($n=9$; 3.2\%), \texttt{.gz} ($n=6$; 2.1\%), and uncompressedformats ($n=4$; 1.4\%). Compression was gzip in 276 files (97.2\%), withcompression detected via magic-byte inspection in 276 files (97.2\%) and viafile-extension fallback in 4 files (1.4\%). The delimiter was a tab characterin 276 files (97.2\%), a space in 4 files (1.4\%), and could not be determinedin 3 files (1.1\%). Headers were detected successfully in 280 of 284 files(98.6\%), and 280 files (98.6\%) were read without error.

Raw file size varied substantially across the dataset, with a median of318~MB (range 0--3,881~MB; 5th--95th percentile: 4.3--1,784~MB). Row countsshowed similarly large variation: the median physical line count was$1.13 \times 10^{7}$ (range 0--$1.62 \times 10^{8}$), with 150 files (52.8\%)containing more than 10 million variants, 52 files (18.3\%) between 1 and 10million, 75 files (26.4\%) between 100,000 and 1 million, and 6 files (2.1\%)fewer than 100,000 variants. The number of columns in the raw file ranged from0 to 29, with the most common values being 12 columns (100 files; 35.2\%),24 columns (36 files; 12.7\%), and 14 columns (30 files; 10.6\%), confirmingthat GWAS files differed substantially in the amount of supplementary annotationthey carried alongside the core summary statistics.

\paragraph{Schema diversity and standard-column mapping.}Despite sharing a common underlying statistical structure, candidate GWAS filesused 78 distinct raw header signatures across 283 files, which collapsed to73 distinct mapped standard-header signatures after automated columnnormalisation. This indicates that a large proportion of files used source-specific column naming conventions that required explicit mapping before any downstream analysis.

Table~\ref{tab} summarises the mapping rate and most common source column names for each GWASRanker standard column. Core positional andassociation columns were mapped at high rates: the effect allele (A1) wasmapped in 280 files (98.6\%), P-value in 280 files (98.6\%), standard error(SE) in 278 files (97.9\%), other allele (A2) in 279 files (98.2\%), andeffect-allele frequency (EAF) in 272 files (95.8\%). Chromosome andbase-pair position were mapped in 271 (95.4\%) and 275 (96.8\%) files,respectively. In contrast, optional or study-design-specific columns showedmuch lower mapping rates: total sample size (N) was available in only 107files (37.7\%), case count (N\_CASES) in 41 files (14.4\%), control count(N\_CONTROLS) in 37 files (13.0\%), INFO score in 36 files (12.7\%), andminor allele frequency (MAF, distinct from EAF) in only 1 file (0.4\%).The median number of unique standard columns mapped per file was 10(range 0--13; SD~=~1.70), and the median number of non-standard unmappedcolumns was 2 (range 0--13; SD~=~2.89).

\begin{table}[H]
\centering
\caption{\textbf{Standard-column mapping rates and source column names across 284 candidate GWAS files.} For each GWASRanker standard column, the number andpercentage of files in which the column was successfully mapped are reported,together with the most common raw source column names observed.}
\label{tab}
\small
\begin{adjustbox}{max width=\textwidth}
\begin{tabular}{llrrl}
\toprule
\textbf{Standard column} & \textbf{Description} & \textbf{Mapped ($n$)} &\textbf{Mapped (\%)} & \textbf{Most common source column names} \\
\midrule
SNP   & Variant identifier        & 279 & 98.2 & \texttt{variant\_id} (35\%); \texttt{rsid} (22\%) \\
CHR   & Chromosome                & 271 & 95.4 & \texttt{chromosome} (92\%); \texttt{CHR} (3\%) \\
BP    & Base-pair position        & 275 & 96.8 & \texttt{base\_pair\_location} (92\%); \texttt{BP} (2\%); \texttt{POS} (1\%) \\
A1    & Effect allele             & 280 & 98.6 & \texttt{effect\_allele} (90\%); \texttt{A1} (2\%); \texttt{Allele1} (2\%) \\
A2    & Other allele              & 279 & 98.2 & \texttt{other\_allele} (79\%); \texttt{other\_allele|alt} (13\%); \texttt{A2} (2\%) \\
BETA  & Effect size (log-OR/beta) & 202 & 71.1 & \texttt{beta} (68\%); \texttt{Effect} (1\%); \texttt{BETA} (1\%) \\
OR    & Odds ratio                & 165 & 58.1 & \texttt{odds\_ratio} (56\%); \texttt{OR} (1\%) \\
Z     & Z-score                   &  17 &  6.0 & \texttt{z\_score} (3\%); \texttt{z} (2\%) \\
SE    & Standard error            & 278 & 97.9 & \texttt{standard\_error} (91\%); \texttt{se} (2\%); \texttt{SE} (2\%) \\
P     & P-value                   & 280 & 98.6 & \texttt{p\_value} (89\%); \texttt{P} (3\%); \texttt{p} (1\%) \\
N     & Total sample size         & 107 & 37.7 & \texttt{n} (30\%); \texttt{all\_total} (5\%); \texttt{N} (2\%) \\
N\_CASES    & Case count          &  41 & 14.4 & \texttt{num\_cases} (13\%); \texttt{N\_case} (1\%) \\
N\_CONTROLS & Control count       &  37 & 13.0 & \texttt{num\_controls} (13\%) \\
EAF   & Effect-allele frequency   & 272 & 95.8 & \texttt{effect\_allele\_frequency} (86\%); \texttt{freq\_effect\_allele} (6\%) \\
MAF   & Minor-allele frequency    &   1 &  0.4 & \texttt{MAF} ($<$1\%) \\
INFO  & Imputation quality score  &  36 & 12.7 & \texttt{INFO} (6\%); \texttt{info\_all} (6\%) \\
DIRECTION & Direction of effect   &  49 & 17.3 & \texttt{direction} (14\%); \texttt{Direction} (2\%) \\
\bottomrule
\end{tabular}
\end{adjustbox}
\end{table}

\paragraph{Effect-size representation and P-value format.}Effect-size representation varied across files. Of the 283 files with availableinformation, 202 (71.1\%) provided a BETA column, 65 (22.9\%) provided an oddsratio requiring log transformation, 9 (3.2\%) provided only a Z-score, and 7(2.5\%) were missing any effect column entirely. P-values were most commonlyreported in standard decimal format (268 files; 94.4\%), with 11 files (3.9\%)using scientific notation, 3 files (1.1\%) missing a P-value column, and 1 file(0.4\%) using a negative-log$_{10}$ transformation that required back-conversion.

\paragraph{Variant identifier and coordinate formats.}Chromosome values were reported in standard numeric format (1--22) in 268 files(94.4\%), were missing in 12 files (4.2\%), and used mixed or non-standardrepresentations in 3 files (1.1\%). Base-pair positions were integers in 268files (94.4\%), missing in 8 files (2.8\%), and non-numeric in 7 files (2.5\%).SNP identifier format showed greater heterogeneity: 138 files (48.6\%) usedrsIDs, 106 files (37.3\%) used mixed or non-standard identifiers, 35 files(12.3\%) used chromosome--position--allele composite identifiers, and 4 files(1.4\%) lacked any identifier column. Allele encoding was predominantlysingle-base uppercase (A, C, G, T) for both the effect allele (242 files;85.2\%) and the other allele (241 files; 84.9\%), with lowercase encoding in 20files each (7.0\%) and mixed or invalid encoding in 18 files each (6.3\%).

\paragraph{Per-file missingness and duplication.}Per-file missingness rates differed substantially across GWAS fields and acrossfiles (Table~\ref{tab}). Allele columns and P-value columnsshowed the lowest missingness: 260 of 283 files had zero missing A1 values, and276 files had zero missing P-values. In contrast, SNP identifier missingnesswas zero in only 175 files, with 92 files showing greater than 10\% missingSNP identifiers (maximum 100\%, indicating files with no usable rsID column).Effect-size missingness exceeded 10\% in 55 files, reflecting the subset offiles that carried only OR or Z-score columns rather than BETA. DuplicatedSNP identifier rates showed the greatest spread: only 88 files had zeroduplicated SNP records, 90 files showed more than 10\% duplication (medianacross all files: 0.54\%), and some files showed 100\% duplication ratesattributable to non-unique composite identifier structures. Palindromic (ambiguous) allele percentage was the most pervasivequality-flag field, with a median of 14.3\% and more than 10\% ambiguousalleles in 251 of 283 files --- an expected consequence of the A/T and C/Gallele pairs that cannot be strand-resolved without allele-frequencyinformation.

\begin{table}[H]
\centering
\caption{\textbf{Per-file missingness and duplication distributions across 283 candidate GWAS files.} For each GWAS field, the number of files with zeromissingness or duplication, fewer than 1\%, and greater than 10\% are reported,together with the median and maximum percentage across all files.}
\label{tab}
\small
\begin{tabular}{lrrrll}
\toprule
\textbf{Field} & \textbf{Files = 0\%} & \textbf{Files $<$ 1\%} &\textbf{Files $>$ 10\%} & \textbf{Median (\%)} & \textbf{Max (\%)} \\
\midrule
SNP identifier      & 175 & 182 &  92 & 0.000 & 100.0 \\
Effect allele (A1)  & 260 & 279 &   4 & 0.000 & 100.0 \\
Other allele (A2)   & 259 & 278 &   5 & 0.000 & 100.0 \\
P-value             & 276 & 279 &   3 & 0.000 & 100.0 \\
Effect size (B/Z)   & 225 & 227 &  55 & 0.000 & 100.0 \\
Duplicated SNP      &  88 & 173 &  90 & 0.542 & 100.0 \\
Ambiguous alleles   &  25 &  28 & 251 & 14.321 & 100.0 \\
Invalid alleles     & 263 & 264 &  17 & 0.000 & 100.0 \\
\bottomrule
\end{tabular}
\end{table}

\paragraph{PRS compatibility and preprocessing requirements.}Despite the heterogeneity described above, 276 of 284 files (97.2\%) containedthe minimum column set required for PRS construction (SNP, A1, BETA or OR, P),and 276 files (97.2\%) were successfully normalised. Seven files (2.5\%) couldnot be normalised safely and were excluded from downstream GWASRanker rankingand PRS analyses. Column renaming was required in 276 files (97.2\%), reflectingthe near-universal use of source-specific rather than GWASRanker-standard columnnames. Decompression or extraction was required in the same 276 files (97.2\%).Effect-size conversion (log-transformation of OR to BETA) was required in 65files (22.9\%). Compatibility with the GWAS-SSF harmonised schema was observedin 204 files (71.8\%).

Among non-standard columns that were present but not mappable to any GWASRanker standard field, the most prevalent were confidence-interval columns(\texttt{ci\_lower} and \texttt{ci\_upper}, each present in 161 files; 56.7\%),GWAS Catalog harmonisation metadata columns (\texttt{hm\_coordinate\_conversion}and \texttt{hm\_code}, each present in 142 files; 50.0\%), andheterogeneity statistics from meta-analyses (\texttt{i2}, \texttt{r2}, and\texttt{q\_pval}, each present in 37--38 files; 13.0--13.4\%). These columnswere retained in the raw files but excluded from GWASRanker feature extractionand PRS construction.

Taken together, Feature1 and Feature2 demonstrated that candidate GWAS summary-statistic files were not directly interchangeable and requiredsystematic file auditing, compression handling, delimiter detection, columnmapping, effect-size standardisation, and schema normalisation before GWASRanker ranking and PRS construction could proceed.

\subsection{Harmonisation and FinalGWAS quality}

Harmonisation and FinalGWAS quality were assessed using Feature3 and Feature5, which captured whether each raw GWAS file could be converted into a standardised \texttt{FinalGWAS.csv} representation and whether the resulting harmonised file was suitable for downstream GWASRanker ranking and PRS analyses. Variation was observed across genome-build detection, GWASLab loading and quality-control outcomes, inferred file format, duplicate variant rates, FinalGWAS scale, allele composition, association signal strength, variant frequency spectrum, sample-size reporting, and downstream tool compatibility.

\paragraph{Genome-build detection.}Of the 283 files with available build information, 179 (63.0\%) were assigned to GRCh38/hg38 and 52 (18.3\%) to GRCh37/hg19, while 52 files (18.3\%) yielded an unknown build that could not be resolved from the available file metadata. Build inference succeeded for 231 of 284 files (81.3\%). Among the 239 files that were successfully loaded by GWASLab, hg38 build-match counts were substantially higher than hg19 match counts for the majority of files (median hg38 match count $1.09 \times 10^{6}$ versus median hg19 match count $1.76 \times 10^{4}$), consistent with the predominance of GRCh38 submissions in the dataset. The build-match ratio varied across a wide range (min $= 0$; max $= 0.774$; median $= 0.017$), indicating that some files had a clear build signal while others had limited coordinate overlap with either reference. Despite this heterogeneity in raw-file build inference, 278 of 284 FinalGWAS files (97.9\%) were ultimately processed under GRCh38, reflecting the pipeline's use of build-lifting where necessary.

\paragraph{GWASLab loading and quality-control removal.}GWASLab loaded 239 of 284 files successfully (84.2\%); the remaining 45 files (15.8\%) failed to load due to structural or encoding issues unresolvable after schema mapping. Among the successfully loaded files, chromosome values were parsed as numeric in 225 files (79.2\%), base-pair position in 238 files (83.8\%), and P-value in 237 files (83.5\%). BETA and SE columns were parsed as numeric in 144 (50.7\%) and 189 (66.5\%) files respectively, with the remainder carrying OR or Z-score columns that required separate handling. The loaded variant count ranged from $2.54 \times 10^{4}$ to $1.62 \times 10^{8}$ (median $1.20 \times 10^{7}$). GWASLab quality-control removal rates were low for the majority of files: 133 files (46.8\%) lost fewer than 1\% of variants, and 66 files (23.2\%) lost between 1\% and 5\%. Ten files (3.5\%) lost between 5\% and 20\% of variants, and 3 files (1.1\%) lost more than 20\%, with a maximum removal rate of 97.2\% in one file. The median QC-removed percentage across all 239 files was 0.41\% (mean 2.33\%; SD 10.76\%), indicating that most files required only modest variant-level filtering but that a tail of files underwent substantially heavier QC attrition.

\paragraph{GWASLab inferred file format.}Among the 239 successfully loaded files, GWASLab assigned the top inferred format as \texttt{ssf} (GWAS-SSF schema) in 150 files (52.8\%) and as \texttt{gwascatalog} in 83 files (29.2\%), with \texttt{pheweb} (3 files; 1.1\%), \texttt{gwaslab} (2 files; 0.7\%), and \texttt{plink} (1 file; 0.4\%) accounting for the remainder. The top format confidence score ranged from 14.73 to 48.38 (median 42.94; mean 39.67), indicating that most files were assigned to their top format with high confidence, but a subset showed lower discrimination between competing format assignments.

\paragraph{Duplicate variants before and after harmonisation.}Duplicate variant records were present at substantial levels in some files before GWASLab quality control. Duplicated SNP identifier counts ranged from 0 to $2.58 \times 10^{7}$ (median $1.33 \times 10^{4}$), duplicated chromosome--position records from 0 to $7.33 \times 10^{6}$ (median $4.26 \times 10^{4}$), and duplicated chromosome--position--allele records from 0 to $1.24 \times 10^{6}$ (median 34). After harmonisation, the FinalGWAS files retained residual duplication: the median duplicated SNP identifier rate was 0.25\% (max 7.33\%), the median duplicated chromosome--position rate was 0.67\% (max 100\%, reflecting files with entirely position-based rather than rsID-based identifiers), and the median duplicated chromosome--position--allele count was 40 (max $1.24 \times 10^{6}$). These results confirm that deduplication and coordinate validation were necessary preprocessing steps for all downstream GWASRanker analyses.

\paragraph{FinalGWAS scale and standard-column completeness.}The 278 files that produced a valid FinalGWAS varied substantially in size and variant count (Table~\ref{tab}). FinalGWAS file sizes ranged from 0.9~MB to 12,978~MB (median 1,047~MB), and loaded variant counts ranged from $1.0 \times 10^{4}$ to $1.62 \times 10^{8}$ (median $1.12 \times 10^{7}$). A total of 147 files (51.8\%) contained more than 10 million variants, 43 files (15.1\%) between 2 and 10 million, 30 files (10.6\%) between 500,000 and 2 million, and 58 files (20.4\%) fewer than 500,000 variants. All 22 autosomes were present in 245 of 284 files (86.3\%). Core association columns were recovered at high rates: SNPID, CHR, POS, BETA, OR, and SE were each present in 278 FinalGWAS files (97.9\%), EA and NEA in 277 files (97.5\%), P in 277 files (97.5\%), and EAF in 269 files (94.7\%). Total sample size (N) was available in 260 files (91.5\%), case count (N\_CASES) in 207 files (72.9\%), control count (N\_CONTROLS) in 206 files (72.5\%), and INFO score in only 58 files (20.4\%).

\begin{table}[H]
\centering
\caption{\textbf{Standard-column availability in 278 FinalGWAS files.} Columns are listed in order of GWASRanker priority. N refers to total sample size; N\_CASES and N\_CONTROLS to case and control counts; EAF to effect-allele frequency; INFO to imputation quality score.}
\label{tab}
\small
\begin{tabular}{llrr}
\toprule
\textbf{Standard column} & \textbf{Description} & \textbf{Present ($n$)} & \textbf{Present (\%)} \\
\midrule
SNPID       & Variant identifier          & 278 & 97.9 \\
CHR         & Chromosome                  & 278 & 97.9 \\
POS         & Base-pair position          & 278 & 97.9 \\
EA          & Effect allele               & 277 & 97.5 \\
NEA         & Non-effect allele           & 277 & 97.5 \\
BETA        & Effect size                 & 278 & 97.9 \\
OR          & Odds ratio                  & 278 & 97.9 \\
SE          & Standard error              & 278 & 97.9 \\
P           & P-value                     & 277 & 97.5 \\
N           & Total sample size           & 260 & 91.5 \\
N\_CASES    & Case count                  & 207 & 72.9 \\
N\_CONTROLS & Control count               & 206 & 72.5 \\
EAF         & Effect-allele frequency     & 269 & 94.7 \\
INFO        & Imputation quality score    &  58 & 20.4 \\
Z           & Z-score                     &   9 &  3.2 \\
MAF         & Minor-allele frequency      &   1 &  0.4 \\
\bottomrule
\end{tabular}
\end{table}

\paragraph{Allele composition and variant types.}The allele composition of FinalGWAS files was dominated by single-base SNPs (median 93.9\%; mean 86.9\%), with indels present at a median of 6.1\% (mean 12.4\%) and multi-character alleles at a median of 3.1\% (mean 9.2\%). Both the indel and multi-character allele distributions were right-skewed, with maximum values of 100\% in a small number of files that contained only indel or multi-allelic records. Palindromic SNPs (A/T or C/G pairs) were present in a median of 15.2\% of variants per file (mean 13.4\%; max 16.3\%), consistent with the expected proportion of strand-ambiguous variants in genome-wide SNP data. Non-ACGT allele percentages had a median of zero but a mean of 6.2\%, indicating that a subset of files carried symbolic or non-standard allele encodings requiring additional filtering. Symbolic alleles were absent from all files.

\paragraph{Association signal and P-value validity.}Association signal strength varied markedly across the 278 FinalGWAS files. The minimum P-value per file ranged from 0 to $3.4 \times 10^{-5}$ (median $9.6 \times 10^{-15}$), reflecting large differences in the strength of the strongest associations present. Genome-wide significant variant counts (P $\leq 5 \times 10^{-8}$) ranged from 0 to 228,500 (median 51), and 71 files (25.0\%) contained no genome-wide significant variants at all. At the other extreme, 129 files (45.4\%) contained 100 or more genome-wide significant variants, and 36 files (12.7\%) contained between 1 and 9. Median P-values across files were centred near the expected uniform-null value (median 0.483; range 0.001--0.902), with the lower end of this range indicating files with genuine genome-wide signal enrichment. P--Z concordance was high overall (median 100\%; mean 94.2\%), but 19 files showed concordance below 80\%, suggesting partial inconsistency between reported P-values and effect sizes that would affect downstream LDSC and PRS analyses.

\paragraph{Variant frequency spectrum.}The variant frequency spectrum differed substantially across files, reflecting differences in imputation panels, array density, and study design. Rare variants (MAF $< 1\%$) constituted a median of 49.5\% of variants per file (mean 44.1\%), low-frequency variants (1--5\%) a median of 8.7\% (mean 11.3\%), and common variants (MAF $> 5\%$) a median of 23.2\% (mean 25.6\%). The median effect-allele frequency across all variants within each file ranged from 0 to 0.964 (median across files: 0.007), indicating that many files were enriched for low-frequency or rare variants, which has direct implications for PRS variant selection and clumping thresholds.

\paragraph{Sample-size reporting.}Total sample size (N) was reported in 260 of 284 files (91.5\%), with 211 files (74.3\%) reporting a constant N across all variants and 49 files (17.3\%) reporting variant-level N. Median total N across files ranged from 681 to 1,713,200 (median across files: 368,600; mean: 336,800). Case counts were available in 207 files, with a median of 22,060 cases (range 69--441,300), and control counts in 206 files, with a median of 297,600 controls (range 608--1,272,000). These sample-size differences are expected to contribute substantially to variation in downstream PRS performance.

\paragraph{Effect-size validity and downstream tool compatibility.}All 278 FinalGWAS files with available effect-size information were stored with a BETA column (97.9\%), reflecting the pipeline's conversion of OR to log-OR during harmonisation. Extreme effect-size values were uncommon: the median extreme OR percentage was 0\% (max 6.1\%), and the median extreme $|\text{BETA}| > 10$ percentage was also 0\% (max 4.0\%). For downstream tool compatibility, 276 FinalGWAS files (97.2\%) met the minimum column requirements for both PRS construction and PLINK, 259 files (91.2\%) met the LDSC minimum requirements, and 268 files (94.4\%) satisfied the GWAS-SSF mandatory column specification. Taken together, Feature3 and Feature5 demonstrated that harmonisation produced a common file structure and effect-size representation across the candidate GWAS files, but did not eliminate biological, statistical, and variant-quality differences between them. These differences were therefore retained as informative GWASRanker features.

\subsection{GWAS-intrinsic statistical quality}

GWAS-intrinsic statistical quality was assessed using Feature4 and Feature6, which captured statistical properties of each harmonised GWAS file independently of target-genotype matching or downstream PRS performance. Variation was observed across sample-size reporting, variant identifier and allele quality, association signal strength, INFO and allele-frequency availability, statistical inflation, beta and standard-error distributions, heritability estimates, chromosomal coverage, and the effect--frequency relationship.

\paragraph{Sample-size reporting.}Total sample size (N) was available in 260 of 284 files (91.5\%). Among files with N information, the reporting style was constant across all variants in 211 files (74.3\%), variant-level in 49 files (17.3\%), and unavailable in 18 files (6.3\%). Median total N per file ranged from 681 to 1,713,000 (median across files: 368,600; mean: 336,800; SD: 225,000), reflecting substantial differences in GWAS power across the candidate dataset. Case counts were available in 207 files (72.9\%) and control counts in 206 files (72.5\%). The N variability ratio (max/min N within a file) had a median of zero, indicating that most files with N information reported it as constant across variants, but a subset showed within-file variation up to a ratio of 1.0, indicating variant-level missingness or cohort-level heterogeneity.

\paragraph{Variant identifier and allele quality.}SNP identifiers were available in rsID format in 253 of 278 files (89.1\%), with 25 files (8.8\%) lacking a usable identifier column. Palindromic SNPs (A/T or C/G pairs) constituted a median of 15.2\% of variants per file (mean 13.4\%; max 16.3\%), consistent with expected genome-wide proportions. Indel percentages varied widely (median 6.1\%; mean 12.4\%; max 100\%), as did multi-character allele percentages (median 3.1\%; mean 9.2\%; max 100\%), with the extreme values indicating a small number of files composed entirely of indel or multi-allelic records. Missing variant identifier rates showed extreme heterogeneity: the median was near zero (0.00045\%) but the mean was 34.2\% and the maximum was 100\%, reflecting files in which rsID-based identifiers were entirely absent. Duplicated variant identifier rates were lower overall (median 0.25\%; max 7.3\%). Allele case was uniformly non-mixed across all 278 files, and lowercase alleles were present in A1 or A2 in 8 files each (2.8\%).

\paragraph{Association signal and P-value distribution.}A P-value column was present in 277 of 284 files (97.5\%). Minimum P-values per file ranged from 0 to $3.4 \times 10^{-5}$ (median $9.6 \times 10^{-15}$), while median P-values were centred near the expected null value of 0.5 (median 0.483; range 0.001--0.902). Genome-wide significant variant counts (P $\leq 5 \times 10^{-8}$) ranged from 0 to 228,500 (median 51; mean 7,755). A total of 71 files (25.0\%) contained no genome-wide significant variants, 36 files (12.7\%) contained 1--9, 42 files (14.8\%) contained 10--99, and 129 files (45.4\%) contained 100 or more. Suggestive association counts (P $\leq 1 \times 10^{-5}$) showed similarly wide variation (median 946; max 490,400). The number of independent genome-wide significant loci (window-based) ranged from 0 to 2,230 (median 6; mean 89): 69 files (24.3\%) had no independent loci, 90 files (31.7\%) had 1--10, 41 files (14.4\%) had 11--50, 38 files (13.4\%) had 51--200, and 29 files (10.2\%) had more than 200 independent loci.

\paragraph{INFO score and allele-frequency availability.}An INFO score column was present in only 58 of 284 files (20.4\%), with INFO missing in a median of 100\% of variants across all files (mean missingness 85.6\%), confirming that imputation quality information was available only for a minority of GWAS files. Among the 40 files with a non-missing median INFO score, values ranged from 0.771 to 1.0 (median 0.898), and fewer than 1\% of variants had INFO below 0.3 in these files. An effect-allele frequency (EAF) column was present in 269 files (94.7\%), while a MAF column was present in only 1 file (0.4\%). The variant frequency spectrum varied substantially: rare variants (MAF $< 1\%$) constituted a median of 49.5\% of variants per file (SD 39.8\%), low-frequency variants (1--5\%) a median of 8.7\% (SD 12.8\%), and common variants ($> 5\%$) a median of 23.2\% (SD 28.2\%), indicating that some files were imputed to a dense panel including rare variants while others reported only common variants.

\paragraph{Statistical inflation.}Genomic-control lambda ($\lambda_{\text{GC}}$) varied from 0.033 to 24.6 across 268 files (median 1.081; mean 1.184; SD 1.471), and mean chi-square ranged from 0.264 to 12.75 (median 1.057). The QQ inflation pattern was classified as mixed in 151 files (53.2\%), uniform inflation with possible confounding in 104 files (36.6\%), and strong top-driven signal in 13 files (4.6\%). Z-score consistency failure exceeded 10\% in 24 files (8.5\%), and the inconsistent Z-score percentage varied from 0 to 95.3\% (median 0.003\%; mean 4.5\%), indicating that a subset of files had systematic discordance between their reported P-values and effect-size estimates.

\paragraph{Beta and standard-error distributions.}Beta distributions showed a wide range of shapes across files. Beta means were near zero for most files but showed extreme outliers in both directions (min $-9.8 \times 10^{8}$; max $1.1 \times 10^{4}$), reflecting scale differences between BETA and OR-derived log-OR files. Beta skewness ranged from $-1{,}752$ to $1{,}860$ and kurtosis from $-0.55$ to $3.7 \times 10^{6}$, indicating that some files contained extreme outlier effect estimates. The beta distribution flag was classified as normal in 169 files (59.5\%), highly skewed in 45 files (15.8\%), large-beta-SD (possible log-OR scale issue) in 25 files (8.8\%), extreme kurtosis in 16 files (5.6\%), and insufficient valid beta in 13 files (4.6\%). Beta outlier percentages at $|z| > 5$ had a median of 0.18\% (max 1.69\%) and at $|z| > 10$ a median of 0.001\% (max 0.21\%). Standard-error distributions showed comparable heterogeneity: SE means ranged from 0.002 to $1.6 \times 10^{5}$, and the SE coefficient of variation ranged from 0.10 to 3,274 (median 0.91). The SE bimodal flag was true in 131 files (46.1\%), indicating that almost half of the files showed evidence of multi-cohort or multi-array SE structure. The SE distribution was classified as normal in 118 files (41.5\%), wide IQR in 101 files (35.6\%), high CV suggesting mixed cohorts in 35 files (12.3\%), and insufficient valid SE in 13 files (4.6\%).

\paragraph{Heritability estimates and LDSC metrics.}Heritability estimation was attempted and succeeded for 225 of 284 files (79.2\%), with 43 files (15.1\%) failing due to LDSC input requirements not being met. Heritability-related features showed the largest variation of any feature group in the GWASRanker dataset, spanning negative estimates, near-zero values, and large positive values across the 225 files for which LDSC completed successfully (Table~\ref{tab}). LDSC $h^2$ estimates ranged from $-0.182$ to $23.65$ (median $0.036$; mean $0.412$; SD $2.37$), with extreme positive values indicating LDSC estimation artefacts in files with very large sample sizes or non-standard effect-size scales. Among the 225 files, 19 (6.7\%) yielded negative $h^2$ estimates, 52 (18.3\%) were near-zero ($h^2 < 0.02$), 93 (32.7\%) were low ($0.02$--$0.10$), 43 (15.1\%) were moderate ($0.10$--$0.30$), and 18 (6.3\%) were high ($h^2 > 0.30$). The LDSC intercept ranged from 0.428 to 2.378 (median 1.032), with 62 files (21.8\%) showing deflation (intercept $< 1.00$), 64 files (22.5\%) near-nominal (1.00--1.05), 39 files (13.7\%) moderate inflation (1.05--1.10), and 60 files (21.1\%) notable inflation ($\geq 1.10$). The sample-size-weighted heritability signal ($N \times h^2$) ranged from $-70{,}590$ to $252{,}000$ (median $8{,}844$): 19 files had negative $N \times h^2$, 12 (4.2\%) had values below 1,000, 88 (31.0\%) between 1,000 and 10,000, 69 (24.3\%) between 10,000 and 50,000, and 37 (13.0\%) above 50,000. LDSC variant removal during cleaning ranged from 0.0001\% to 100\% (median 10.5\%; mean 45.1\%): 41 files (14.4\%) lost fewer than 5\% of variants, 109 files (38.4\%) lost 5--20\%, and 116 files (40.8\%) lost more than 50\%, reflecting the large proportion of files in which LDSC-incompatible variant identifiers or missing rsIDs required extensive cleaning before heritability estimation.

\begin{table}[H]
\centering
\caption{\textbf{Heritability-related feature variation across 284 candidate GWAS files.} LDSC estimation succeeded for 225 files; remaining fields are reported for all files with available data. $N \times h^2$ is the sample-size-weighted heritability signal used by GWASRanker as a composite GWAS-quality indicator. The LDSC ratio reflects the proportion of inflation attributable to confounding rather than true polygenic signal; values near zero indicate signal-driven inflation and values near 1 indicate confounding-driven inflation.}
\label{tab}
\small
\begin{adjustbox}{max width=\textwidth}
\begin{tabular}{lrrrrrrrr}
\toprule
\textbf{Metric} & \textbf{$n$} & \textbf{Missing (\%)} & \textbf{Min} & \textbf{p5} & \textbf{Median} & \textbf{Mean} & \textbf{p95} & \textbf{Max} \\
\midrule
LDSC $h^2$                    & 225 & 20.8 & $-0.182$ & $-0.021$ & $0.036$ & $0.412$ & $0.464$ & $23.65$ \\
LDSC $h^2$ SE                 & 225 & 20.8 & $0.001$ & $0.001$ & $0.007$ & $0.069$ & $0.268$ & $1.781$ \\
LDSC intercept                & 225 & 20.8 & $0.428$ & $0.762$ & $1.032$ & $1.067$ & $1.487$ & $2.378$ \\
LDSC intercept SE             & 225 & 20.8 & $0.005$ & $0.006$ & $0.013$ & $0.080$ & $0.352$ & $0.798$ \\
LDSC ratio                    & 161 & 43.3 & $0.006$ & $0.082$ & $0.239$ & $0.683$ & $2.771$ & $13.96$ \\
LDSC ratio SE                 & 161 & 43.3 & $0.008$ & $0.009$ & $0.050$ & $0.637$ & $4.049$ & $9.183$ \\
LDSC mean $\chi^2$            & 225 & 20.8 & $0.901$ & $1.006$ & $1.194$ & $1.428$ & $2.989$ & $5.327$ \\
LDSC $\lambda_{\text{GC}}$    & 225 & 20.8 & $0.895$ & $0.993$ & $1.163$ & $1.287$ & $2.067$ & $4.236$ \\
$N \times h^2$                & 225 & 20.8 & $-70{,}590$ & $-7{,}550$ & $8{,}844$ & $28{,}470$ & $133{,}600$ & $252{,}000$ \\
LDSC variant removal (\%)     & 268 &  5.6 & $0.000$ & $0.028$ & $10.45$ & $45.10$ & $100.0$ & $100.0$ \\
\bottomrule
\end{tabular}
\end{adjustbox}
\end{table}

\paragraph{Chromosomal coverage and variant density.}All 22 autosomes were represented in 245 of 284 files (86.3\%). Median variant density across autosomes ranged from 0 to 57,920 variants per Mb (median across files: 4,066 per Mb; mean 5,344 per Mb), and minimum per-chromosome variant density ranged from 0 to 44,480 per Mb (median 2,860 per Mb). The percentage of zero-coverage 1~Mb genomic bins ranged from 3.8\% to 85.7\% (median 5.9\%; mean 8.8\%), indicating that most files had good autosomal coverage but a subset had substantial genomic gaps, consistent with array-based rather than imputed GWAS data. Chr22 variant counts ranged from 0 to 2,262,000 (median 145,300).

\paragraph{Effect--frequency relationship.}The absolute beta--log(MAF) Spearman correlation, which captures the expected inverse relationship between effect size and minor allele frequency, ranged from $-0.947$ to $0.301$ across 205 files (median $-0.628$; mean $-0.564$). The negative median value is consistent with the expected genetic architecture in which rarer variants tend to have larger effect sizes. However, 63 files (22.2\%) lacked sufficient valid beta--MAF pairs for this calculation, and 205 files (72.2\%) yielded a descriptive annotation rather than a quality-control failure, indicating that the metric was informative but not universally available. Taken together, Feature4 and Feature6 demonstrated that candidate GWAS files differed substantially in intrinsic statistical signal, calibration, allele composition, frequency spectrum, heritability, and chromosomal coverage, and that these differences were not removed by harmonisation. They were therefore retained as informative GWASRanker features for model training.
\paragraph{Per-phenotype heritability variation.}Heritability estimates varied substantially not only across the full candidate GWAS dataset but also within each phenotype, reflecting the heterogeneity of study designs, sample sizes, and analytical approaches represented among the candidate GWAS files for each trait. Across all 225 files for which LDSC heritability estimation succeeded, $h^2$ ranged from $-0.182$ to $23.65$ (median $0.036$), $N \times h^2$ from $-70{,}590$ to $252{,}000$ (median $8{,}844$), and the LDSC intercept from $0.428$ to $2.378$ (median $1.032$). Table~\ref{tab} summarises the per-phenotype heritability, inflation, and association signal variation across all 13 phenotypes represented in the GWASRanker candidate dataset.

\begin{table}[H]
\centering
\caption{\textbf{Per-phenotype heritability, inflation, and association signal variation across 284 candidate GWAS files.} For each phenotype, the number of candidate GWAS files, the number for which LDSC $h^2$ estimation succeeded, the median total sample size, the $h^2$ range and median, the $N \times h^2$ median, the median genomic-control $\lambda_{\text{GC}}$, the median number of independent genome-wide significant loci, and the best (minimum) P-value observed across all files are reported. $h^2$ values marked with $\dagger$ include one or more LDSC estimation artefacts attributable to extreme sample sizes or non-standard effect-size scales. GORD = gastro-oesophageal reflux disease. IBS = irritable bowel syndrome. BMI = body mass index.}
\label{tab}
\small
\begin{adjustbox}{max width=\textwidth}
\begin{tabular}{lrrrlrrrr}
\toprule
\textbf{Phenotype} &\textbf{Files} &\textbf{$h^2$ succeeded} &\textbf{Median $N$} &\textbf{$h^2$ range (median)} &\textbf{$N \times h^2$ median} &\textbf{$\lambda_{\text{GC}}$ median} &\textbf{Indep.\ loci median} &\textbf{Best $P$} \\
\midrule
Asthma &34 & 32 & 375,460 &$-0.022$--$23.6^{\dagger}$ (0.052) &12,979 & 1.14 & 35 &$2.3 \times 10^{-156}$ \\
\addlinespace
Blood pressure medication &4 & 3 & 198,160 &$0.141$--$3.06^{\dagger}$ (2.70) &203,890 & 1.05 & 3 &$2.6 \times 10^{-16}$ \\
\addlinespace
BMI &40 & 25 & 399,130 &$0.118$--$12.7^{\dagger}$ (0.225) &90,652 & 1.25 & 398 &$< 10^{-300}$ \\
\addlinespace
Cholesterol-lowering medication &4 & 3 & 192,570 &$0.038$--$0.826$ (0.051) &16,209 & 1.04 & 41 &$< 10^{-300}$ \\
\addlinespace
Depression &32 & 27 & 297,630 &$-0.027$--$0.118$ (0.024) &4,200 & 1.10 & 1 &$8.5 \times 10^{-25}$ \\
\addlinespace
GORD / gastric reflux &13 & 10 & 381,660 &$-0.099$--$0.071$ ($-0.015$) &$-4{,}333$ & 1.08 & 0 &$1.1 \times 10^{-22}$ \\
\addlinespace
Hayfever / allergic rhinitis &10 & 8 & 192,840 &$-0.012$--$0.100$ (0.022) &3,640 & 1.07 & 1 &$4.5 \times 10^{-21}$ \\
\addlinespace
High cholesterol &5 & 5 & 394,630 &$0.001$--$0.403$ (0.039) &16,728 & 1.15 & 89 &$< 10^{-300}$ \\
\addlinespace
Hypertension &28 & 18 & 331,750 &$-0.182$--$0.760$ (0.082) &26,434 & 1.06 & 6 &$< 10^{-300}$ \\
\addlinespace
Hypothyroidism &26 & 22 & 417,290 &$-0.042$--$0.570$ (0.030) &13,102 & 1.05 & 22 &$< 10^{-300}$ \\
\addlinespace
IBS &26 & 22 & 330,400 &$0.000$--$0.129$ (0.023) &2,900 & 1.01 & 0 &$3.9 \times 10^{-10}$ \\
\addlinespace
Migraine &14 & 14 & 397,560 &$-0.032$--$0.085$ (0.019) &8,007 & 1.05 & 8 &$1.7 \times 10^{-40}$ \\
\addlinespace
Osteoarthritis &48 & 36 & 401,500 &$-0.062$--$0.241$ (0.026) &7,751 & 1.07 & 5 &$2.7 \times 10^{-45}$ \\
\bottomrule
\end{tabular}
\end{adjustbox}
\end{table}

Within individual phenotypes, the range of $h^2$ estimates was similarly wide. For asthma, $h^2$ ranged from $-0.022$ to $23.6$ (median $0.052$) with $N \times h^2$ spanning $-1{,}220$ to $227{,}470$, the extreme upper value reflecting an LDSC artefact in one large-sample file. For body mass index, $h^2$ ranged from $0.118$ to $12.7$ (median $0.225$) and $N \times h^2$ from $6{,}993$ to $157{,}300$ (median $90{,}652$), the highest effective signal of any phenotype, consistent with the large well-powered GWAS available for this trait. For hypertension, $h^2$ ranged from $-0.182$ to $0.760$ (median $0.082$) with the LDSC intercept spanning $0.777$ to $1.910$, indicating substantial heterogeneity in both true signal and inflation across candidate files. For depression, $h^2$ ranged from $-0.027$ to $0.118$ (median $0.024$) with $\lambda_{\text{GC}}$ ranging from $0.033$ to $24.6$, the widest inflation range observed across any phenotype, attributable to one severely inflated file in this group. Phenotypes with consistently low $h^2$ included gastro-oesophageal reflux (median $h^2 = -0.015$; five of ten files with negative estimates; median $N \times h^2 = -4{,}333$), irritable bowel syndrome (median $h^2 = 0.023$; all estimates near-zero or low; median $N \times h^2 = 2{,}900$), and hayfever or allergic rhinitis (median $h^2 = 0.022$; one negative estimate). Phenotypes with moderate to high $h^2$ in at least some files included high cholesterol (median $h^2 = 0.039$; one file reaching $h^2 = 0.403$; median $N \times h^2 = 16{,}728$), hypothyroidism (median $h^2 = 0.030$; two files with $h^2 > 0.30$; $N \times h^2$ ranging from $-16{,}104$ to $189{,}010$), cholesterol-lowering medication ($h^2$ range $0.038$--$0.826$), and blood pressure medication ($h^2$ range $0.141$--$3.06$; median $N \times h^2 = 203{,}890$, reflecting large-sample proxy-phenotype GWAS). Association signal strength mirrored heritability rankings: body mass index had the highest median independent locus count (median 398; max 2,230; best $P < 10^{-300}$), followed by cholesterol-lowering medication (median 41 loci), high cholesterol (median 89 loci), asthma (median 35 loci; best $P = 2.3 \times 10^{-156}$), and hypothyroidism (median 22 loci), while gastro-oesophageal reflux (median 0 loci; best $P = 1.1 \times 10^{-22}$), irritable bowel syndrome (median 0 loci; best $P = 3.9 \times 10^{-10}$), and hayfever (median 1 locus; best $P = 4.5 \times 10^{-21}$) showed the weakest signal. These within-phenotype and between-phenotype differences in heritability, inflation, and signal strength confirm that automated GWAS ranking is necessary even when all candidate files nominally target the same trait, and that $N \times h^2$ and independent locus count are particularly discriminating features for GWASRanker model training.

\subsection{Reference-panel and heritability features}

Reference-panel integration was assessed across two feature sets: Feature~6, which capturesallele-frequency concordance between each GWAS file and a matched reference panel, andFeature~7, which quantifies variant overlap against HapMap3, dbSNP151, 1000~Genomes~EUR,and a PLINK BIM reference, and records whether LD-based clumping was attempted.

\subsubsection*{Feature~6: reference allele-frequency alignment}

Each of the 284 GWAS files was matched to one of two build-specific allele-frequencyreference panels derived from the 1000~Genomes~HapMap3 resource: the hg38 panel wasselected for 169 files (59.5\%) and the hg19 panel for 99 files (34.9\%), with theremaining 16 files (5.6\%) receiving no reference assignment. Despite successful panelselection, reference AF correlation and mean absolute AF difference could not be computedfor any file; the HapMap3 match-method field recorded \texttt{failed\_reference\_empty}for all 268 files that carried a reference annotation. This indicates that thereference-overlap step ran but returned an empty result set, most likely because theGWAS variant identifiers and the reference panel coordinates could not be joined at thetime of processing. Consequently, quantitative AF-concordance metrics (correlation,mean absolute difference, overlap count) are entirely absent from this dataset and cannotcontribute to ranking or quality stratification.

\subsubsection*{Feature~7: external reference-panel overlap}

Feature~7 extended the overlap assessment to four external references. Summary statisticsfor the reference panel sizes and the GWAS-to-reference overlap counts are presented inTable~\ref{tab}.

\begin{table}[h]
\centering
\caption{Feature~7 reference-panel sizes and GWAS overlap counts across 278 files withnon-missing Feature~7 data (out of 284 total).}
\label{tab}
\begin{tabular}{lrrrr}
\toprule
Reference & Panel rsIDs & Panel chr & GWAS rsID overlap & GWAS chr overlap \\
\midrule
HapMap3 (PAN)     & 0 (all files)       & $1.20\times10^{6}$ & 0 & 0 \\
dbSNP151          & $6.4$--$7.9\times10^{7}$ & $6.7$--$7.9\times10^{7}$ & 0 & 0 \\
1000G EUR         & 0 (all files)       & $7.0$--$8.5\times10^{7}$ & 0 & 0 \\
PLINK BIM (hg19)  & $1.67\times10^{6}$  & $1.67\times10^{6}$ & 0 & 0 \\
\bottomrule
\end{tabular}
\end{table}

Across all four references, both rsID-based and chromosome-position-based overlap countswere zero for every file. The HapMap3 rsID panel size was also recorded as zero,suggesting that the rsID index for that reference was not populated during the Feature~7run. The dbSNP151 chr panel sizes were consistent and large (median$6.7\times10^{7}$), confirming that the reference itself was loaded correctly; theabsence of any GWAS overlap therefore reflects a failure in the join rather than a missingreference file. The 1000~Genomes~EUR chr panel was similarly populated (median$7.0\times10^{7}$), yet yielded no overlapping variants. These results indicate asystematic upstream issue, most plausibly that the FinalGWAS files passed to Feature~7contained no usable variant coordinates at the time of processing, consistent with the\texttt{missing\_columns} error recorded for 97.9\% of files in the PLINK clumpinginput step.

\subsubsection*{Feature~7: LD-based clumping}

PLINK-based LD clumping was not successfully attempted for any file. The clumping inputstep reported \texttt{missing\_columns} for 278 of 284 files (97.9\%), indicating thatthe pipeline could not locate a valid P-value column in the clumping input. Of the 284files, only 51 (18.0\%) had a compatible hg19 PLINK bfile available(\texttt{1000G\_phase3\_common\_norel}); the remaining 227 files (79.9\%) were flagged\texttt{no\_plink\_bfile\_for\_hg38}, as no hg38 PLINK bfile was provided. Where PLINK2was nominally available (51 files), it returned exit code~7 for all of them, confirmingthat clumping did not complete. The genome-build distribution recorded by Feature~7matched the Feature~3 and Feature~6 assignments: 227 files (79.9\%) were processed ashg38 and 51 files (18.0\%) as hg19, with the build sourced from a prior feature record(\texttt{prior\_feature}) in all 278 cases.

\subsubsection*{Summary}

The reference-panel and heritability integration features (Feature~6 and Feature~7) werestructurally complete in that reference files were selected and the pipeline modulesexecuted, but they produced no quantitative overlap or concordance metrics. AF correlation,HapMap3 overlap, dbSNP151 overlap, 1000~Genomes~EUR overlap, and PLINK BIM overlap wereall zero or missing across the full cohort of 284 GWAS files. LD-based clumping was notperformed for any file. These features therefore contribute no discriminatory signal tothe current ranking analysis, and any scoring dimension that relies on reference-panelconcordance or clump-based locus counts must be treated as unavailable for this dataset.

\subsection{Target-genotype compatibility}

Target-genotype compatibility was assessed across Feature~8 and the relevant subset ofFeature~9. Feature~8 quantifies how well each candidate GWAS file overlaps with a fixedtarget genotype dataset in terms of variant identity, allele concordance, andallele-frequency agreement. Feature~9 summarises the final count of PRS-ready variantsthat survive all upstream filters and are present in the target.

\subsubsection*{Target genotype dataset}

The target genotype dataset was identical across all 270 files for which Feature~8completed successfully (out of 284 total; 14 files had no Feature~8 data). The target BIMcontained exactly 619,700 variants on the hg38/GRCh38 reference assembly, all withchromosome-position coordinates and rsIDs available (99.34\% rsID coverage). The targetcarried no duplicate rsIDs or duplicate chromosome-position entries. All 619,700 variantshad MAF available, with a median target MAF of 0.072, indicating a panel of commonvariants. The target comprised 586 samples (median; range 582--586), of whomapproximately 442 were female and 143 male. Of the 270 files with phenotype information,232 (81.7\%) had a binary phenotype encoded in PLINK 1/2 case-control format, with amedian of 75 cases and 511 controls per file and a median effective sample size of 262.The remaining 38 files (13.4\%) were classified as continuous or non-standard, allcorresponding to the BMI phenotype (mean 27.33, SD 4.98).

\subsubsection*{GWAS variant scale and identifier availability}

GWAS files varied enormously in scale, from 10,000 to 161.9~million variants, with amedian of 11.4~million and a 90th-percentile range of approximately 0.18--32~million(Table~\ref{tab}). Chr coordinates were available for virtuallyall variants in most files (median 100\%, mean 98.8\%), and allele information wassimilarly near-complete (median 100\%, mean 99.3\%). rsID availability was more variable:the median was 100\% but the mean was only 67.8\%, with a standard deviation of 45.4\%,indicating that a substantial subset of files carried no rsIDs at all. Palindromic(A/T and C/G) SNPs constituted a median of 14.3\% of variants per file (range0--16.0\%). Duplicate rsID rates were low (median 0.15\%), but duplicate chr rateswere highly variable, with a median of 0.34\% but a mean of 6.3\% and a p95 of 89.9\%,indicating that a small number of files had near-complete chr non-uniqueness.

\begin{table}[h]
\centering
\caption{GWAS variant count and identifier availability across 270 Feature~8-completefiles.}
\label{tab}
\begin{tabular}{lrrrrr}
\toprule
Metric & Min & p5 & Median & p95 & Max \\
\midrule
Total variant count      & 10,000 & 184,200 & 11,440,000 & 31,970,000 & 161,900,000 \\
rsID available (\%)      & 0      & 0       & 100        & 100        & 100         \\
Chr available (\%)   & 0      & 99.95   & 100        & 100        & 100 \\
Allele available (\%)    & 0      & 99.96   & 100        & 100        & 100         \\
MAF available (\%)       & 0      & 0       & 100        & 100        & 100         \\
Palindromic SNP (\%)     & 0      & 0       & 14.3       & 15.3       & 16.0 \\
Duplicate rsID (\%)      & 0      & 0       & 0.15       & 2.0        & 3.9 \\
Duplicate chr (\%)   & 0      & 0       & 0.34       & 89.9       & 89.9        \\
\bottomrule
\end{tabular}
\end{table}

\subsubsection*{GWAS-target variant overlap}

Overlap between each GWAS and the target BIM was computed by both rsID andchromosome-position matching. rsID-based overlap yielded a median of 581,600 variants(94.5\% of the target, but only 4.1\% of the GWAS), reflecting the fact that the targetpanel is small relative to most GWAS files. Chr overlap was substantially lowerin absolute terms (median 9,056 variants; 1.5\% of target), suggesting that rsID matchingwas the primary join strategy and that chr alone recovered far fewer variants whenrsIDs were unavailable or inconsistent. The bimodal distributions of both overlap metrics— with near-zero p5 values alongside high medians — indicate that a substantial minorityof files achieved little or no overlap, corresponding to the files with zero rsIDavailability noted above.

\subsubsection*{Allele matching}

Among rsID-overlapping variants, allele matching decomposed as follows: direct strandmatch (median 306,700), strand-flipped match (median 73,200), palindromic ambiguous(median 34,730), allele mismatch (median 12), and missing allele (median 0). Theresulting rsID-based allele-compatible count had a median of 546,900, representing93.8\% of rsID-overlapping variants. The small mismatch and missing-allele countsindicate high allele concordance where rsID overlap was achieved. Chr allelematching showed a similar pattern but lower absolute counts (median compatible count7,028; 74.9\% of chr overlap), with proportionally more mismatches (median 857),consistent with the less specific nature of chr matching.

\subsubsection*{Allele-frequency concordance}

MAF concordance between the GWAS and the target was assessed separately for rsID-matchedand chr variants. For rsID-based matching, MAF correlation was high (median0.996, mean 0.965; available for 200 of 284 files), with a mean absolute MAF differenceof 0.007 (median), indicating excellent allele-frequency agreement where rsID overlap wasachieved. Chr MAF correlation was markedly lower (median 0.718, mean 0.637;available for 191 files) with a median absolute difference of 0.012 and a mean of 0.029,reflecting the inclusion of non-uniquely matched variants when falling back to positionaljoining. These results confirm that rsID-based matching is the more reliable approach forallele-frequency harmonisation in this dataset.

\begin{table}[h]
\centering
\caption{GWAS-target MAF concordance by matching strategy across files with availableMAF data.}
\label{tab}
\begin{tabular}{lrrrr}
\toprule
Strategy & n files & MAF correlation (median) & MAF correlation (mean) & Mean abs.\ diff.\ (median) \\
\midrule
rsID-based      & 200 & 0.996 & 0.965 & 0.007 \\
Chr   & 191 & 0.718 & 0.637 & 0.012 \\
\bottomrule
\end{tabular}
\end{table}

\subsubsection*{Final usable variants and PRS readiness}

After applying all overlap and allele-compatibility filters, the median final usablevariant count was 555,700 per file (89.7\% of target BIM; 2.8\% of GWAS), with a p95 of575,700 and a maximum of 579,800. The near-ceiling recovery relative to the target BIMin the upper half of files demonstrates that, when rsID overlap is achieved, the targetpanel is well-covered. The PRS-ready and in-target count from Feature~9 was slightlylower (median 534,500; 2.1\% of GWAS; 100\% of ready variants for 91.1\% of files),with the small reduction relative to Feature~8 reflecting additional PRS-specificreadiness filters applied upstream. Files with zero final usable variants correspondedto those with no rsID or chr overlap, which in turn reflects the zero-rsID-availableGWAS files identified above. The complete summary of final usable and PRS-ready variantcounts is given in Table~\ref{tab}.

\begin{table}[h]
\centering
\caption{Final usable and PRS-ready target-overlapping variant counts across 270Feature~8-complete files.}
\label{tab}
\begin{tabular}{lrrrrr}
\toprule
Metric & Min & p5 & Median & p95 & Max \\
\midrule
Final usable variant count (Feature~8)            & 0 & 0 & 555,700 & 575,700 & 579,800 \\
Final usable \% of target BIM                     & 0 & 0 & 89.7    & 92.9    & 93.6 \\
Final usable \% of GWAS                           & 0 & 0 & 2.8     & 9.4     & 14.1    \\
PRS-ready and in-target count (Feature~9)         & 0 & 0 & 534,500 & 577,100 & 583,100 \\
PRS-ready \% of GWAS (Feature~9)                  & 0 & 0 & 2.1     & 8.3     & 15.4    \\
PRS-ready \% of ready variants (Feature~9)        & 0 & 0 & 100     & 100     & 100     \\
\bottomrule
\end{tabular}
\end{table}

\subsection{PRS readiness, clumping, and pruning}

\subsubsection*{PRS readiness}

PRS readiness was evaluated across seven sequential criteria: SNPID availability, validP-value, effect size, alleles, clump-readiness, score-readiness, and full PRS readiness.Of the 270 files with Feature~9 data (14 of 284 had no Feature~9 output), P-values werenear-universally available (median 100\%, mean 99.6\%), alleles were similarly complete(median 100\%, mean 99.3\%), and effect sizes were available for essentially all variantsin most files (median 100\%, mean 98.1\%). SNPID availability was the most variablecriterion, with a mean of 67.8\% and a standard deviation of 45.4\%, indicating thatroughly one-third of files carried no rsIDs. No beta values were derived from ORconversion (all files had direct beta columns). The fully PRS-ready variant count had amedian of 534,500 and a mean of 319,400, with the bimodal distribution reflecting thesplit between files with rsID-based target overlap (median $\sim$534K) and those without(zero). Of 270 files, 24 (8.5\%) had zero fully-ready variants, 80 (28.2\%) had between1 and 100K, 28 (9.9\%) had between 100K and 500K, and 138 (48.6\%) had between 500K and2M. Of the ready variants, the median fraction that were also present in the target was100\% (mean 91.1\%), confirming that once PRS readiness was achieved the target overlapwas largely complete.

Duplicate SNPIDs were present at a low median rate of 0.15\% (mean 0.32\%, max 3.9\%),and the clump-ready count (median 545,800) was slightly higher than the score-ready andfull-ready count (median 534,500), with the difference attributable to variants thatpassed clumping filters but failed final score-readiness checks.

\subsubsection*{LD clumping}

Clumping was performed using PLINK with a fixed window of 250~kb and $r^{2} = 0.1$,across seven P-value thresholds from the genome-wide significance threshold($P \leq 5\times10^{-8}$) to the null ($P \leq 1.0$). Clumping succeeded for all 270files at every threshold. The clump input comprised a median of 555,700 variants(2.8\% of the full GWAS), reflecting the target-overlap filter applied upstream.Table~\ref{tab} summarises clump counts across all files at each threshold,and Table~\ref{tab} gives per-phenotype medians, minima, and maxima atthe genome-wide significance threshold and at $P \leq 1\times10^{-3}$.

At the genome-wide significance threshold, 103 of 270 files (36.3\%) produced zero clumps,while 167 files (58.8\%) yielded at least one clump (median 33, p95 921, max 3,117).Files with non-zero GW clumps spanned a median of 10 chromosomes. The mean SP2 count(proxy variants per index SNP) at this threshold was 19.1 (median 11.0, max 315),indicating moderate LD block sizes. As the P-value threshold was relaxed, the proportionof files with zero clumps fell monotonically: 23.9\% at $P \leq 10^{-6}$, 15.5\% at$P \leq 10^{-5}$, 10.9\% at $P \leq 10^{-4}$, 7.4\% at $P \leq 10^{-3}$, and 7.0\% atboth $P \leq 0.05$ and $P \leq 1.0$. The 20 persistently zero-clump files correspond tothose with no usable target-overlapping variants. At $P \leq 1.0$, files with non-zeroclumps had a median of 205,500 independent loci (p95 210,900), spanning all 22 autosomesin nearly every case.

\begin{table}[h]
\centering
\caption{LD clump count distribution across 270 Feature~9-complete files at eachP-value threshold. Clumping parameters: window 250~kb, $r^{2} = 0.1$.}
\label{tab}
\begin{tabular}{lrrrrrrr}
\toprule
Threshold & n zero & n $>$ 0 & Min & p5 & Median & p95 & Max \\
\midrule
$P \leq 5\times10^{-8}$ & 103 & 167 & 0 & 0    & 3      & 572    & 3,117 \\
$P \leq 1\times10^{-6}$ & 68  & 202 & 0 & 0    & 6      & 901    & 4,925 \\
$P \leq 1\times10^{-5}$ & 44  & 226 & 0 & 0    & 21     & 1,369  & 7,511 \\
$P \leq 1\times10^{-4}$ & 31  & 239 & 0 & 0    & 98     & 2,287  & 12,602 \\
$P \leq 1\times10^{-3}$ & 21  & 249 & 0 & 0    & 641    & 5,311  & 22,678 \\
$P \leq 0.05$           & 20  & 250 & 0 & 0    & 20,270 & 39,990 & 75,086 \\
$P \leq 1.0$            & 20  & 250 & 0 & 0    & 201,200& 210,900& 211,700\\
\bottomrule
\end{tabular}
\end{table}

\begin{table}[h]
\centering
\caption{Per-phenotype clump counts at $P \leq 5\times10^{-8}$ and$P \leq 1\times10^{-3}$, and LD-pruned variant counts. Values are median(min--max) across files within each phenotype.}
\label{tab}
\resizebox{\textwidth}{!}{%
\begin{tabular}{lrrrrrr}
\toprule
Phenotype & Files &\multicolumn{2}{c}{Clumps at $P \leq 5\times10^{-8}$} &\multicolumn{2}{c}{Clumps at $P \leq 1\times10^{-3}$} &Pruned-in (median) \\
\cmidrule(lr){3-4}
\cmidrule(lr){5-6}
& & Median & Min--Max & Median & Min--Max & \\
\midrule
Asthma                        & 34 & 49    & 0--361    & 878    & 0--3,050   & 358,822 \\
Blood pressure medication     &  4 &  0.5  & 0--2      &   5    & 0--20      &     758 \\
BMI                           & 40 & 115   & 0--3,117  & 2,323  & 0--22,678  & 357,348 \\
Cholesterol-lowering medic.   &  4 &  65   & 0--208    &   523  & 0--2,938   & 182,246 \\
Depression                    & 32 &  0    & 0--15     &   560  & 0--1,394   & 360,008 \\
GORD / gastric reflux         & 13 &  0    & 0--62     &     4  & 0--1,259   &   1,200 \\
Hayfever / allergic rhinitis  & 10 &  0    & 0--5      &     5  & 0--14      &   4,264 \\
High cholesterol              &  5 & 150   & 126--322  & 1,662  & 904--3,578 & 362,907 \\
Hypertension                  & 28 &  2    & 0--387    &   663  & 0--5,400   & 352,683 \\
Hypothyroidism / myxoedema    & 26 &  23   & 0--397    &   595  & 0--3,688   & 362,371 \\
Irritable bowel syndrome      & 26 &  0    & 0--4      &   454  & 0--918     & 322,888 \\
Migraine                      & 14 &  3    & 0--14     &   672  & 0--1,159   & 356,756 \\
Osteoarthritis                & 48 &  4.5  & 0--168    &   719  & 0--3,094   & 360,200 \\
\bottomrule
\end{tabular}}
\end{table}

The phenotypes with the strongest genome-wide signal were BMI (median 115 GW clumps,max 3,117), high cholesterol (median 150), and asthma (median 49), consistent with thelarge discovery sample sizes for these traits. In contrast, blood pressure medication,hayfever, and GORD had near-zero median GW clumps, reflecting both smaller discoveryGWAS and smaller target-overlap variant pools. At the lenient $P \leq 10^{-3}$ threshold,the rank order was broadly preserved: BMI had the highest median (2,323), while GORD(median 4) and hayfever (median 5) remained near-zero.

\subsubsection*{LD pruning}

LD pruning was performed with a 200~kb sliding window, step size 50, and $r^{2} = 0.25$.Pruning succeeded for 250 of 270 files (88.0\%); the remaining 20 (7.0\%) were skippeddue to an empty extract file, corresponding to the same zero-overlap files identifiedabove. Among files where pruning ran, the median prune-in count was 353,200 variants(63.7\% of the input), and the median prune-out count was 200,200 (36.0\% of input),indicating that approximately one-third of target-overlapping variants were in LD withanother retained variant at this threshold. The prune-in count was highly consistentacross files in the upper half of the distribution (p95 364,800, max 365,700), reflectingthe fixed target BIM size. Per-phenotype pruned-in medians are given inTable~\ref{tab}; GORD and hayfever were clear outliers with far fewerretained variants (medians 1,200 and 4,264 respectively), driven by their smallerrsID-based target overlap rather than any difference in LD structure.

\subsection{Source GWAS and target Fisher GWAS concordance}

Target-side GWAS association testing was performed using either a Fisher exact test(binary phenotypes) or a linear regression (continuous phenotypes) on the target cohort,providing an independent estimate of variant-level association for concordance assessment.Of the 270 Feature~9 files, 250 produced an association result file. The associationmethod was linear regression for 36 files (all BMI) and Fisher for 1 file; the remainderused Fisher by default for binary traits. Fisher input comprised a median of 555,700variants, of which a median of 559,600 were common to both the GWAS and targetassociation results. Allele compatibility between the source GWAS and target Fisherresults was high: a median of 100\% of common variants were allele-compatible (mean79.3\%), with direct strand match the dominant category (median 326,500), followed bystrand-flipped match (median 82,980), and negligible mismatch or missing-allele counts(both median 0).

Despite good allele concordance, the concordance between source GWAS effect sizes andtarget Fisher statistics was very low across all metrics. The beta correlation betweensource GWAS and target Fisher betas was near-zero for all common variants (median 0.004,mean 0.006, range $-$0.058--0.061) and for allele-compatible variants only (median 0.006,mean 0.008). The absolute beta correlation was similarly low (median 0.175, mean 0.129).The mean absolute beta difference between source GWAS and target was 0.316 (median), andP-value correlations were effectively zero (median 0.002 for all common variants). Theseresults are expected and appropriate: the target cohort comprised only 586 individuals,giving negligible power to replicate genome-wide association signals from GWAS of tens tohundreds of thousands of participants. The near-zero concordance therefore does notindicate data quality problems but rather the inherent power differential betweendiscovery and target datasets. The Fisher/linear association results are used internally for allele orientation checking and not as a primary quality metric.

\subsection{PRS performance outcomes}

PRS performance was evaluated using both PLINK C+T and PRSice-2, with cross-validatedAUC (binary phenotypes) or $R^{2}$ (continuous phenotypes) as the primary metric acrossa five-fold cross-validation design. Of the 284 candidate GWAS files, 246 (86.6\%)produced a successful best-result file for both tools. Of the remaining 38 files, 38were flagged as missing for PLINK and 32 as failed for PRSice-2, corresponding to thefiles with zero PRS-ready target-overlapping variants identified inSection~\ref{sec}. All 246 successful files used five folds; one PLINK fileused a single fold. The performance metric was AUC for 210 files (73.9\%, binaryphenotypes) and $R^{2}$ for 36 files (12.7\%, BMI as the sole continuous phenotype).

\subsubsection{PRS performance: PLINK C+T}
\label{sec}

\paragraph{Model performance distributions.}Table~\ref{tab} summarises the distribution of PLINK test-set performanceacross all 246 successful files for the pure PRS, null, and best combined models.The pure PRS test AUC had a median of 0.549 (mean 0.528), with 58 files (20.4\%)falling below chance ($\mathrm{AUC} < 0.50$) and 25 files (8.8\%) achieving AUC$\geq 0.80$. The null model (covariates only) showed a markedly bimodal distribution:for binary phenotypes the null model was driven by the high intercept of the targetcohort phenotype prevalence, yielding a median test null AUC of 0.580 (mean 0.436),but for BMI (the continuous phenotype, $n=36$ files) the null model $R^{2}$ wassystematically negative (median $-0.669$, mean $-0.668$), reflecting the absence ofuseful covariates for a continuous trait in this target cohort. The best combined model(PRS plus covariates) had a median test AUC of 0.599 (mean 0.490), with 36 files(12.7\%) below chance (all BMI) and 38 files (13.4\%) achieving AUC $\geq 0.80$.

\begin{table}[h]
\centering
\caption{PLINK C+T test-set performance distribution across 246 successful files.AUC is reported for binary phenotypes; $R^{2}$ for BMI. Values below chance for BMIreflect the negative $R^{2}$ of a continuous trait with no covariate adjustment.}
\label{tab}
\resizebox{\textwidth}{!}{%
\begin{tabular}{lrrrrrrrr}
\toprule
Model & n & Min & p5 & p25 & Median & p75 & p95 & Max \\
\midrule
Test pure PRS   & 246 & $-$0.037 & $-$0.029 & 0.505 & 0.549 & 0.593 & 0.885 & 0.976 \\
Train pure PRS  & 246 & 0.001    & 0.012    & 0.516 & 0.550 & 0.590 & 0.883 & 0.979 \\
Test null model & 246 & $-$0.748 & $-$0.679 & 0.528 & 0.580 & 0.645 & 0.804 & 0.956 \\
Train null model& 246 & 0.574    & 0.576    & 0.843 & 0.865 & 0.913 & 0.936 & 0.997 \\
Test best model & 246 & $-$0.738 & $-$0.635 & 0.550 & 0.599 & 0.686 & 0.898 & 0.968 \\
Train best model& 246 & 0.577    & 0.585    & 0.849 & 0.882 & 0.934 & 0.989 & 1.000 \\
\bottomrule
\end{tabular}}
\end{table}

\paragraph{Incremental PRS contribution.}The incremental AUC of the best model over the null (best minus null on the test set)had a median of 0.011 (mean 0.055), indicating that the PRS contribution abovecovariates was modest for most files. Of 246 files, 25 (8.8\%) showed no gain($\Delta\mathrm{AUC} \leq 0$), 91 (32.0\%) showed marginal gain (0--0.01), 72 (25.4\%)modest gain (0.01--0.05), 18 (6.3\%) moderate gain (0.05--0.10), and 40 (14.1\%)substantial gain ($\geq 0.10$). The incremental AUC of best over pure PRS had a medianof 0.031 and a mean of $-$0.037, with the negative mean driven by BMI files where thenull model performs poorly and the pure PRS already captures most variance.

\paragraph{Generalisation gap.}The pure PRS generalisation gap (train minus test) was small and well-controlled, witha median of 0.008 (mean 0.015, max 0.097); only 10 files (3.5\%) exceeded a gap of0.05. The best model generalisation gap was substantially larger (median 0.277, mean0.373, max 1.315), with 244 files (85.9\%) exceeding 0.05. This is expected andarises because the training null model AUC is very high (median 0.865, reflectingcross-validated overfitting to case-control prevalence in the small target cohort of586 individuals), while the test null AUC is much lower. The large best-model gaptherefore reflects the covariate component rather than PRS overfitting.

\paragraph{Best P-value threshold.}The optimal PLINK C+T P-value threshold was highly variable across files. The mostcommon single threshold was $P \leq 1.0$ (68 files, 23.9\%), meaning the fullunthresholded PRS performed best for those GWAS. The remaining files spanned thresholdsfrom $3.4 \times 10^{-10}$ to 1.0, with a median of 0.026 and a mean of 0.324,indicating that moderately lenient thresholds were generally optimal.

\subsubsection{PRS performance: PRSice-2}
\label{sec}

\paragraph{Model performance distributions.}PRSice-2 performance was closely comparable to PLINK across all metrics (Table~\ref{tab}). The test pure PRS AUC had a median of 0.548 (mean 0.520),with 69 files (24.3\%) below chance and 26 (9.2\%) achieving AUC $\geq 0.80$. The testbest model AUC had a median of 0.599 (mean 0.501), with 36 files (12.7\%) below chanceand 43 (15.1\%) at AUC $\geq 0.80$. As with PLINK, the null model for BMI wassystematically negative (median $-0.644$, constant across files since PRSice-2 uses asingle null model per phenotype).

\begin{table}[h]
\centering
\caption{PRSice-2 test-set performance distribution across 246 successful files.}
\label{tab}
\resizebox{\textwidth}{!}{%
\begin{tabular}{lrrrrrrrr}
\toprule
Model & n & Min & p5 & p25 & Median & p75 & p95 & Max \\
\midrule
Test pure PRS   & 246 & $-$0.043 & $-$0.029 & 0.484 & 0.548 & 0.594 & 0.894 & 0.980 \\
Train pure PRS  & 246 & 0.000    & 0.017    & 0.517 & 0.550 & 0.592 & 0.889 & 0.981 \\
Test null model & 246 & $-$0.644 & $-$0.644 & 0.528 & 0.588 & 0.662 & 0.795 & 0.947 \\
Train null model& 246 & 0.570    & 0.570    & 0.835 & 0.857 & 0.904 & 0.931 & 0.996 \\
Test best model & 246 & $-$0.660 & $-$0.621 & 0.539 & 0.599 & 0.699 & 0.912 & 0.974 \\
Train best model& 246 & 0.571    & 0.577    & 0.842 & 0.877 & 0.927 & 0.987 & 0.998 \\
\bottomrule
\end{tabular}}
\end{table}

\paragraph{Incremental PRS contribution.}The incremental AUC (best minus null) had a median of 0.007 (mean 0.054), slightlylower than PLINK. Of 246 files, 50 (17.6\%) showed no gain, 93 (32.7\%) marginal gain,48 (16.9\%) modest gain, 9 (3.2\%) moderate gain, and 46 (16.2\%) substantial gain.The pure PRS generalisation gap was small (median 0.005, mean 0.022), with 30 files(10.6\%) exceeding 0.05. The best model generalisation gap was again large (median 0.252,mean 0.358), for the same reasons as PLINK.

\paragraph{Best P-value threshold.}PRSice-2 searched a finer grid of thresholds. The most common optimal threshold was$P \leq 10^{-5}$ (39 files, 13.7\%), followed by $P \leq 0.05$ (33 files, 11.6\%) and$P \leq 0.10$ (20 files, 7.0\%). The median optimal threshold was 0.25 (mean 0.333),broadly consistent with PLINK but with a finer resolution of optimal values between$10^{-5}$ and 1.0.

\paragraph{PLINK vs PRSice-2 comparison.}In paired comparisons across 246 files, the median difference in test best-modelperformance (PLINK minus PRSice-2) was $-$0.009 (mean $-$0.010), indicating PRSice-2was marginally superior on average. PRSice-2 outperformed PLINK in 173 of 246 files(70.3\%), PLINK outperformed PRSice-2 in 73 files (29.7\%), and no file was tied. Themaximum advantage for either tool was modest ($\leq 0.17$ AUC units), confirming thatthe two methods produced closely comparable rankings for most GWAS files.

\paragraph{Per-phenotype performance.}Table~\ref{tab} summarises test best-model AUC and incremental AUC byphenotype for both tools. The highest absolute AUC was achieved for blood pressuremedication and cholesterol-lowering medication (PLINK median 0.921 and 0.942respectively), reflecting the strong phenotypic signal of medication-use proxies in thetarget cohort rather than GWAS quality per se. Hypertension (median 0.802) and highcholesterol (median 0.914) also showed consistently high AUC. In contrast, BMI producedsystematically negative best-model performance (PLINK median $-$0.577, PRSice-2 median$-$0.593), a consequence of the negative null $R^{2}$ for continuous traits in thisanalysis rather than a failure of the GWAS files themselves. Among the remaining binaryphenotypes, depression (median 0.582), asthma (median 0.544), hayfever (median 0.539),and osteoarthritis (median 0.586) showed modest AUC values, while hypothyroidism(median 0.643), GORD (median 0.686), IBS (median 0.660), and migraine (median 0.611)showed moderate AUC. Within-phenotype variation in test AUC was substantial for mosttraits, with ranges spanning 0.20--0.40 AUC units across GWAS files for the samephenotype, confirming that GWAS file quality rather than phenotype alone drivesperformance differences.

\begin{table}[h]
\centering
\caption{Per-phenotype PRS test-set performance summary for PLINK C+T and PRSice-2. AUC is reported for binary phenotypes; $R^{2}$ for BMI. Incremental AUC is bestmodel minus null model on the test set.}
\label{tab}
\resizebox{\textwidth}{!}{%
\begin{tabular}{lrcccccc}
\toprule
Phenotype & $n$ &\multicolumn{3}{c}{PLINK test best model} &\multicolumn{3}{c}{PRSice-2 test best model} \\
\cmidrule(lr){3-5}
\cmidrule(lr){6-8}& & Min & Median & Max & Min & Median & Max \\
\midrule
Asthma                       & 34 & 0.518 & 0.544 & 0.838 & 0.528 & 0.536 & 0.882 \\
Blood pressure medication    &  4 & 0.920 & 0.921 & 0.952 & 0.938 & 0.940 & 0.940 \\
BMI (continuous, $R^{2}$)   & 40 & $-$0.738 & $-$0.577 & $-$0.238 & $-$0.660 & $-$0.593 & $-$0.235 \\
Cholesterol-lowering medic.  &  4 & 0.939 & 0.942 & 0.944 & 0.947 & 0.947 & 0.949 \\
Depression                   & 32 & 0.549 & 0.582 & 0.891 & 0.591 & 0.596 & 0.922 \\
GORD / gastric reflux        & 13 & 0.659 & 0.686 & 0.827 & 0.693 & 0.698 & 0.808 \\
Hayfever / allergic rhinitis & 10 & 0.502 & 0.539 & 0.578 & 0.523 & 0.532 & 0.548 \\
High cholesterol             &  5 & 0.897 & 0.914 & 0.945 & 0.907 & 0.909 & 0.962 \\
Hypertension                 & 28 & 0.788 & 0.802 & 0.883 & 0.795 & 0.801 & 0.901 \\
Hypothyroidism / myxoedema   & 26 & 0.588 & 0.643 & 0.967 & 0.634 & 0.646 & 0.974 \\
Irritable bowel syndrome     & 26 & 0.629 & 0.660 & 0.784 & 0.662 & 0.666 & 0.805 \\
Migraine                     & 14 & 0.574 & 0.611 & 0.943 & 0.591 & 0.601 & 0.946 \\
Osteoarthritis               & 48 & 0.551 & 0.586 & 0.899 & 0.570 & 0.581 & 0.913 \\
\bottomrule
\end{tabular}}
\end{table}

\section*{Supplementary correlation analysis of GWAS-level features and PLINK pure PRS performance}
\label{sec}

To further investigate which GWAS-level properties were associated with downstream PRS performance, we performed an exploratory correlation analysis between the merged GWAS feature table and PLINK pure PRS performance metrics. This analysis was designed to identify whether variation in GWAS structure, signal strength, sample size, variant completeness, allele quality, target overlap, and PRS-readiness features was associated with PLINK-derived PRS performance.

Three PLINK pure PRS outcomes were evaluated: training pure PRS performance, test pure PRS performance, and the PLINK train--test generalisation gap. Null-model and best-model performance metrics were excluded from this supplementary analysis because the objective was to focus specifically on pure PRS performance. To avoid circularity, all PLINK-derived columns and all PRSice-2-derived columns were removed from the candidate feature set before correlation analysis.

The merged table contained 1015 columns. After excluding PLINK and PRSice-2 derived columns, 966 columns were considered as candidate GWAS-level predictors. Among these, 575 columns were native numerical columns and 391 were string or object columns. String/object columns were converted to numerical values where possible, including Boolean-like values and numeric-formatted values. A total of 106 string/object columns were successfully transformed into numerical values, while 310 string/object columns were excluded as non-numerical. After removing columns with insufficient non-missing observations or no variation, 526 numerical or numerical-convertible GWAS-level features were retained for correlation testing.

Spearman rank correlation was calculated between each retained feature and each of the three PLINK pure PRS outcomes, resulting in 1506 valid feature--outcome correlations. Multiple testing correction was performed using the Benjamini--Hochberg false discovery rate procedure across all feature--outcome tests. Correlations with global FDR-adjusted (q < 0.05) were considered statistically significant. In total, 608 correlations were significant at nominal (p < 0.05), 481 remained significant after global FDR correction, and 519 were significant after target-wise FDR correction.

For PLINK test pure PRS performance, 502 valid feature--outcome correlations were calculated. Among these, 281 were nominally significant and 234 remained significant after global FDR correction. The strongest positive associations were observed for GWAS signal-strength and PRS-readiness features, including the number of clumped genome-wide significant loci, the number of variants passing genome-wide significance, maximum ($-\log_{10}(P)$), and counts or percentages of genome-wide and suggestive associations. These results indicate that GWAS datasets with stronger association signals tended to produce better PLINK pure PRS performance.

Sample-size-related features were also positively associated with PLINK test pure PRS performance. These included median GWAS sample size, maximum GWAS sample size, minimum GWAS sample size, and total sample size estimates. Although these associations were more moderate than those observed for direct GWAS signal-strength features, they support the expected relationship between larger GWAS sample size and improved PRS performance.

Several data-quality features showed negative associations with PLINK test pure PRS performance. These included missing SNP identifier percentage, duplicated SNP percentage, invalid chromosome percentage, non-autosomal variant percentage, indel percentage, multi-character allele percentage, missing effect-size information, missing allele information, and reduced variant retention during quality control. These findings suggest that incomplete or poorly harmonised GWAS files can reduce downstream PRS performance, even when association signals are present.

Target compatibility features were also associated with PLINK performance. Positive associations were observed for GWAS--target rsID overlap, allele-compatible overlap, and available variant counts. This suggests that PRS performance was partly influenced by the extent to which GWAS variants could be matched and harmonised with the target genotype dataset.

Overall, this supplementary analysis indicates that PLINK pure PRS performance was influenced by multiple GWAS-level properties, including association signal strength, sample size, variant completeness, allele representation, target genotype overlap, and data-quality features. The strongest associations were observed for GWAS signal-strength and PRS-readiness features, while missingness, duplication, invalid variant annotations, and allele-formatting issues were generally associated with reduced performance. The complete ranked list of FDR-significant correlations is provided in Supplementary Table~\ref{tab} and the full machine-readable table should be provided as an accompanying supplementary file.

\begin{table}[htbp]
\centering
\caption{Representative FDR-significant correlations between GWAS-level features and PLINK test pure PRS performance. The complete ranked correlation output should be provided as a separate supplementary table or machine-readable file.}
\label{tab}
\small
\begin{tabular}{p{0.24\textwidth} p{0.38\textwidth} r r r}
\toprule
\textbf{Feature group} & \textbf{Representative feature} & \textbf{Spearman ($\rho$)} & \textbf{Global FDR (q)} & \textbf{Direction} \\
\midrule
GWAS signal strength & Number of clumped genome-wide significant loci & 0.740 & ($1.37 \times 10^{-13}$) & Positive \\
GWAS signal strength & Number of clumped loci at ($P \leq 1 \times 10^{-6}$) & 0.700 & ($1.04 \times 10^{-11}$) & Positive \\
GWAS signal strength & Number of clumped loci at ($P \leq 1 \times 10^{-5}$) & 0.673 & ($1.41 \times 10^{-10}$) & Positive \\
GWAS signal strength & Number of variants with ($P \leq 5 \times 10^{-8}$) & 0.667 & ($2.52 \times 10^{-10}$) & Positive \\
GWAS signal strength & Maximum ($-\log_{10}(P)$) & 0.641 & ($2.45 \times 10^{-9}$) & Positive \\
GWAS signal strength & Number of suggestive associations & 0.551 & ($1.39 \times 10^{-15}$) & Positive \\
GWAS signal strength & Number of genome-wide significant associations & 0.548 & ($1.42 \times 10^{-15}$) & Positive \\
GWAS signal strength & Percentage of genome-wide significant associations & 0.504 & ($3.83 \times 10^{-13}$) & Positive \\
Sample size & Maximum GWAS sample size & 0.293 & ($1.55 \times 10^{-4}$) & Positive \\
Sample size & Median GWAS sample size & 0.293 & ($1.55 \times 10^{-4}$) & Positive \\
Sample size & Total GWAS sample size & 0.293 & ($1.55 \times 10^{-4}$) & Positive \\
Sample size & Minimum GWAS sample size & 0.275 & ($4.40 \times 10^{-4}$) & Positive \\
Variant completeness & Missing SNP identifier percentage & -0.410 & ($1.51 \times 10^{-8}$) & Negative \\
Variant completeness & Missing variant identifier percentage & -0.352 & ($1.87 \times 10^{-6}$) & Negative \\
Variant completeness & Missing effect-size information percentage & -0.287 & ($1.78 \times 10^{-4}$) & Negative \\
Variant completeness & Missing allele information percentage & -0.275 & ($3.45 \times 10^{-4}$) & Negative \\
Variant quality & Duplicated SNP percentage & -0.483 & ($5.28 \times 10^{-12}$) & Negative \\
Variant quality & Duplicated chromosome-position percentage & -0.322 & ($1.84 \times 10^{-5}$) & Negative \\
Variant quality & Invalid chromosome percentage & -0.330 & ($1.07 \times 10^{-5}$) & Negative \\
Variant quality & Non-autosomal variant percentage & -0.330 & ($1.07 \times 10^{-5}$) & Negative \\
Variant quality & Indel percentage & -0.311 & ($4.04 \times 10^{-5}$) & Negative \\
Variant quality & Multi-character allele percentage & -0.314 & ($3.27 \times 10^{-5}$) & Negative \\
Target compatibility & GWAS--target rsID overlap count & 0.362 & ($3.80 \times 10^{-3}$) & Positive \\
Target compatibility & Allele-compatible rsID overlap count & 0.356 & ($4.60 \times 10^{-3}$) & Positive \\
Target compatibility & Final GWAS rsID available count & 0.347 & ($6.01 \times 10^{-3}$) & Positive \\
Target compatibility & Final GWAS allele available percentage & 0.323 & ($1.14 \times 10^{-2}$) & Positive \\
Frequency and imputation quality & Median INFO score & 0.485 & ($1.04 \times 10^{-2}$) & Positive \\
Frequency and imputation quality & Mean INFO among common variants & 0.611 & ($7.50 \times 10^{-4}$) & Positive \\
Frequency and imputation quality & Low INFO percentage among common variants & -0.595 & ($1.13 \times 10^{-3}$) & Negative \\
Frequency and imputation quality & Median effect allele frequency & 0.279 & ($1.57 \times 10^{-3}$) & Positive \\
\bottomrule
\end{tabular}
\end{table}

\begin{table}[htbp]
\centering
\caption{Summary of the supplementary PLINK pure PRS correlation analysis.}
\label{tab}
\small
\begin{tabular}{lr}
\toprule
\textbf{Analysis component} & \textbf{Count} \\
\midrule
Total columns in merged table & 1015 \\
PLINK pure PRS outcome columns & 3 \\
Candidate columns after PLINK/PRSice-2 exclusion & 966 \\
Native numerical candidate columns & 575 \\
Native string/object candidate columns & 391 \\
String/object columns converted to numerical values & 106 \\
String/object columns excluded as non-numerical & 310 \\
Final numerical candidate features tested & 526 \\
Valid feature--outcome correlations & 1506 \\
Nominally significant correlations, (p < 0.05) & 608 \\
Globally FDR-significant correlations, (q < 0.05) & 481 \\
Target-wise FDR-significant correlations, (q < 0.05) & 519 \\
\bottomrule
\end{tabular}
\end{table}

\end{document}
